% TOP_PROBLEM: {"id": 30006686, "title": "Basing One-Way Functions on NP-Hardness", "category_id": 15, "category": {"id": 15, "name": "computer_science", "display_name": "Computer Science", "description": "Computational complexity, algorithms, and theoretical CS.", "slug": "computer-science", "order_index": 15, "created_at": "2026-07-31T15:26:25.671Z"}, "status": "open"}
% FIRST_POSTED: 2026-09-27
% !TeX program = lualatex
% ATTEMPT_STATUS: unresolved
\documentclass[11pt]{article}
\usepackage[margin=25mm]{geometry}
\usepackage{fontspec}
\setmainfont{Latin Modern Roman}
\usepackage{amsmath,amssymb,mathtools}
\usepackage{unicode-math}
\setmathfont{Latin Modern Math}
\usepackage{xurl,hyperref}
\hypersetup{colorlinks=true,urlcolor=blue}
\setcounter{secnumdepth}{0}
\setlength{\parindent}{0pt}
\setlength{\parskip}{5pt}
\setlength{\emergencystretch}{3em}
\begin{document}
\section*{\texorpdfstring{Basing One-Way Functions on NP-Hardness}{Basing One-Way Functions on NP-Hardness}}\label{30006686}
\subsection{Definitions and mathematical statement}
A worst-case hardness assertion excludes an efficient algorithm correct on every input; average-case inversion hardness bounds success on $f(U_n)$, where $U_n$ is uniform on $\{0,1\}^n$. The desired construction and reduction would establish
\[
 \mathsf P\ne\mathsf{NP}\quad\Longrightarrow\quad
 \exists\text{ a one-way function family}.
\]
More operationally, an efficient inverter with nonnegligible success would be transformed into an efficient solver for an NP-hard decision problem. The class of allowed reductions matters: an obstruction to a restricted black-box reduction does not by itself refute the unrestricted goal. The catalogue leaves this reduction model broad.
\par\smallskip\textit{Formulation and concept references:} \href{https://api.openalex.org/works/doi:10.1145/1132516.1132614}{[S1]}.

\subsection{Short English statement}
Can worst-case NP-hardness alone be turned into the average-case inversion hardness needed for one-way functions?
\subsection{Research attempt}

\paragraph{Scope, process, and a correction to the source context.}
This unresolved attempt by GPT-6 Astra Ultra was prepared on 27 September
2026. The strategy is to make the missing worst-case-to-average-case step
quantitative, then test it on exact distributions. The assumption is
$\mathsf P\ne\mathsf{NP}$ alone. Replacing it by a conjecture about randomized
algorithms, or by average-case hardness, would change the question.

The author-hosted version of Akavia--Goldreich--Goldwasser--Moshkovitz
\href{https://people.csail.mit.edu/akavia/2006-stocAGGM.pdf}{[E1]}, corresponding to \href{https://api.openalex.org/works/doi:10.1145/1132516.1132614}{[S1]}, studies restricted reductions. Its nonadaptive
result gives a consequence $\mathsf{coNP}\subseteq\mathsf{AM}$ from a reduction
of the stated kind. Crucially, the authors' 2010 erratum \href{https://people.csail.mit.edu/akavia/AGGM_errata.pdf}{[E2]} identifies a
gap in the more general adaptive argument; it preserves the nonadaptive
analysis and an adaptive special case with polynomially bounded fibers.
The unrestricted adaptive assertion from the older abstract is not used.
Even the unaffected result is a conditional obstruction to a reduction
class, not a disproof of the full implication in the catalogue.

The recent references also retain extra hypotheses. The zero-knowledge
result \href{https://arxiv.org/abs/2602.17651}{[S2]} assumes $\mathsf{NP}\not\subseteq\mathsf{ioP/poly}$ and particular
nontrivial zero-knowledge protocols. The boundary-complexity result \href{https://eccc.weizmann.ac.il/report/2025/125/}{[S3]}
states its derandomization assumption for its deterministic-complexity
characterization. Hirahara \href{https://eccc.weizmann.ac.il/report/2023/037/}{[S4]} requires NP-hardness of the specified
meta-computational approximation problem as well as worst-case NP hardness.
None of those additional inputs follows merely by renaming
$\mathsf P\ne\mathsf{NP}$. Their statements were checked for scope; their
full proofs are not independently certified here.

\paragraph{Quantifiers and the easy direction.}
Write $f_n:\{0,1\}^n\to\{0,1\}^{m(n)}$, with polynomially bounded output
length and one uniform polynomial-time evaluation algorithm. For an inverter
$I$, set
\[
 s_I(n)=\Pr_{X\leftarrow\{0,1\}^n,I}
       [I(1^n,f_n(X))\in f_n^{-1}(f_n(X))].
\]
The preimage must have length $n$, but need not be the sampled $X$.
One-wayness requires $s_I$ to be negligible for every probabilistic
polynomial-time $I$. Its negation for this particular family says that
some such $I$ and constant $c>0$ have $s_I(n)>n^{-c}$ on infinitely many
lengths. It does not say that one inverter succeeds almost always, or
succeeds at every length. A proposed reduction must not strengthen this
negation without an argument.

The reverse implication, existence of a one-way function implies
$\mathsf P\ne\mathsf{NP}$, has a short proof. If P equals NP, consider the
language of triples $(1^n,y,u)$ such that some $n$-bit extension $z$ of
prefix $u$ satisfies $f_n(z)=y$. The witness $z$ is polynomial in the
input length, and evaluation and prefix checking are polynomial-time, so
this language would be in P. Start with the empty prefix and ask whether
it can be extended after appending zero. Keep zero if so, and otherwise
keep one. On $y$ in the image, induction retains an extendible prefix.
After $n$ queries the result is a valid preimage. This is one deterministic
polynomial-time inverter on every image value, contradicting one-wayness.
The converse direction requested here is precisely what this search
argument does not supply.

In fact the same construction rules out one-way functions under the weaker
assumption $\mathsf{NP}\subseteq\mathsf{BPP}$. Use the bounded-error decider
for the same prefix language, and reduce each query error to $1/(3n)$ by
majority repetition. For example an odd number of repetitions at least
$27n$ suffices by Chebyshev's inequality, since each trial has bias at
least $1/6$ from one half. Even with adaptively selected prefixes, the
uniform error bound holds conditional on each history. A union bound on
the $n$ queries gives a valid preimage with probability at least $2/3$
on every image value. All repetition costs are polynomial. Thus deriving
one-way functions from the catalogue premise would also imply
$\mathsf P\ne\mathsf{NP}\Longrightarrow\mathsf{NP}\not\subseteq\mathsf{BPP}$;
this additional consequence cannot be silently taken as the premise.

\paragraph{Lemma 33.1: worst-case hardness survives extreme density loss.}
Let $L$ be any NP-complete binary language and define
\begin{equation}\label{r33:padding}
 L_{\rm pad}=\{1^{n^2}0x:\ |x|=n,\ x\in L\}.
\end{equation}
Then $L_{\rm pad}$ is NP-complete, but the deterministic algorithm that
always says no has exponentially small error on uniform strings of each
sufficiently large length.

\emph{Proof.} Mapping $x$ to $1^{|x|^2}0x$ is a polynomial-time many-one
reduction. Membership in the padded language is verified by checking the
format and then using the NP witness for $x\in L$. Thus both hardness and
NP membership hold. At length $M=n^2+n+1$, at most $2^n$ strings have
this format. The proportion accepted by the language, and hence the error
of the constant-no algorithm, is at most
\[
                 2^{n-M}=2^{-n^2-1}.
\]
At every other length the language is empty. For $n\ge1$, $n^2+1\ge M/2$,
so the error bound is at most $2^{-M/2}$ on the nonempty length classes.
\hfill$\square$

This is not a sparse NP-complete set in the usual polynomial-cardinality
sense: the format permits $2^{\Theta(\sqrt M)}$ strings. The argument
therefore does not invoke, or contradict, sparse-hard-set theorems. It
shows directly that deterministic worst-case NP-hardness says nothing
like constant error on the uniform input distribution. A distribution
supported on the padded format would behave differently. Choosing the
right distribution is part of a reduction, not an innocuous encoding
convention.

Likewise, an efficiently generated instance need not be distributed
uniformly among distinct possible instances. If $F$ maps $m$ uniform
seed bits to an instance, its output mass at $y$ is
$2^{-m}|F^{-1}(y)|$. Large fibers are sampled more often. An algorithm
which is good on this distribution may omit a small collection of
particular worst-case targets completely. The following lemma states a
sufficient condition for transferring its guarantee.

\paragraph{Lemma 33.2: a distribution-transfer certificate.}
Fix $n$ and let $\mu_n$ be the law of $f_n(U_n)$. An inverter $I$ has
failure probability
\[
 \epsilon_n=\sum_y\mu_n(y)e_I(y),\qquad
 e_I(y)=\Pr_I[f_n(I(1^n,y))\ne y],
\]
where invalid-length outputs count as failures. Consider a reduction on
an input $x$ that prepares $q$ nonadaptive queries of this same length.
Write $D_{x,j}$ for the marginal law of query $j$. Suppose, for every
$x,j,y$, that
\begin{equation}\label{r33:domination}
                    D_{x,j}(y)\le\kappa\mu_n(y).
\end{equation}
Suppose also that its polynomial-time postprocessing always gives the
correct answer whenever all returned preimages are valid, regardless of
which valid preimages are returned. Then plugging in fresh independent
runs of $I$ makes its error at most $q\kappa\epsilon_n$.

\emph{Proof.} Because the queries are prepared without using inversion
answers, the failure probability at query $j$ is
\[
 \sum_yD_{x,j}(y)e_I(y)
 \le\kappa\sum_y\mu_n(y)e_I(y)=\kappa\epsilon_n.
\]
A union bound over the queries finishes the proof. The queries themselves
may be correlated; independence between their values is not required.
Only their marginals and the freshly run inverter coins enter this bound.
\hfill$\square$

If $q,\kappa,n$ are polynomially bounded in $|x|$, the reduction and inverter
are efficient, and $\epsilon_n\le1/(3q\kappa)$ at all required lengths,
this gives a bounded-error randomized solver. The lemma does not construct
the distributions satisfying \eqref{r33:domination}, the sound postprocessing,
or the required inverter. For adaptive queries, the same proof works if
the domination hypothesis holds after \emph{every possible prior history},
with fresh inverter coins at the next step. An unconditional average over
histories cannot simply be substituted for that conditional hypothesis.
This is an elementary transfer lemma, not a replacement proof for \href{https://people.csail.mit.edu/akavia/2006-stocAGGM.pdf}{[E1]}.

\paragraph{A sharp failure of average-to-worst-case transfer.}
Take the identity function on $n$ bits and let $I$ return its input except
that it fails on $0^n$. Its average failure on $\mu_n=U_n$ is $2^{-n}$.
A query distribution concentrated on $0^n$ nevertheless has failure one.
The least possible domination factor for that query is $\kappa=2^n$,
so the bound from Lemma 33.2 is exactly one. Calling the average error
negligible does not turn it into a useful worst-case bound.

This example uses an easy function deliberately: it refutes a purported
inference about probability distributions without assuming any unproved
hard function. It is not a lower bound against reductions allowed to
exploit the function's code. In particular, the identity has a perfect
inverter; the example concerns a reduction which must be sound when
provided with the specified high-average-success inverter.

The opposite endpoint is equally important. A merely nonnegligible
average success $s_I(n)$ does not supply the small failure
$\epsilon_n$ used by Lemma 33.2. Let $p(y)$ be the success probability
on an individual image value. Repeating $I$ independently $t$ times and
verifying its answers yields average success
\begin{equation}\label{r33:repetition}
       \sum_y\mu_n(y)\bigl(1-(1-p(y))^t\bigr).
\end{equation}
If $p(y)$ equals one on a set of mass $\eta$ and zero elsewhere, this
quantity is $\eta$ for every $t$. Repetition removes bad coins on a good
instance; it cannot create success on instances the inverter never solves.
A worst-case-to-average-case reduction needs an additional mechanism to
redistribute its target into that good region while preserving a way to
recover the original answer.

\paragraph{Two further logical gaps in the proposed reduction.}
First, a randomized solver for an NP-complete language yields
$\mathsf{NP}\subseteq\mathsf{BPP}$. It does not, by definition alone, yield
$\mathsf P=\mathsf{NP}$. The latter requires deterministic deciders.
Thus an argument contradicting only $\mathsf{NP}\not\subseteq\mathsf{BPP}$
would establish a theorem under a different hardness assumption. To use
exactly the catalogue premise, one must also justify the removal of that
randomness, or derive another genuine contradiction to $\mathsf P\ne\mathsf{NP}$.
No such step is supplied here.

Second, a reduction which needs a good inverter at a computable length
$n=n(|x|)$ cannot infer this from nonnegligible success on an unknown
infinite set of lengths. The successful lengths can have arbitrarily
large gaps from the proposed embedding lengths. Padding to a larger
length helps only if an appropriate good length can be found within the
permitted polynomial budget and inversion can be transferred back. Neither
property follows from infinitude alone. These are uniformity and quantifier
issues, not objections to the possibility of a more sophisticated reduction.

\paragraph{Verification, continuation, and outcome.}
The companion exact checks enumerate the padded format at small lengths,
verify its density formula, and compute the identity-function failure
example and \eqref{r33:repetition} with rational arithmetic. They also
check the domination inequality on finite distributions. The written
proofs establish the general lemmas; finite examples do not decide any
complexity-class separation.

\textbf{UNRESOLVED.} The attempt proves an explicit density-loss example
and a quantitative sufficient transfer lemma, and identifies exactly
where their hypotheses fail for a generic inversion claim. A successful
continuation needs a uniform construction and reduction that transports
worst-case NP instances to the actual image distribution, tolerates the
nonnegligible and infinitely-often success given by failure of one-wayness,
and resolves the deterministic-versus-randomized premise. The restricted
literature barriers and these examples do not rule out all such reductions.
No one-way function has been derived from $\mathsf P\ne\mathsf{NP}$ here.

\subsection{Sources}
\begin{itemize}
\item[S1]A. Akavia, O. Goldreich, S. Goldwasser and D. Moshkovitz, 'On basing one-way functions on NP-hardness', Proc. 38th Annual ACM Symp. on Theory of Computing (STOC), 2006, pp. 701-710. DOI:10.1145/1132516.1132614.
\url{https://api.openalex.org/works/doi:10.1145/1132516.1132614}.
\textit{Formulation source.}
\item[S2]Non-Trivial Zero-Knowledge Implies One-Way Functions.
\url{https://arxiv.org/abs/2602.17651}.
\textit{Further reference; not a proof endorsement.}
\item[S3]Hardness Along the Boundary: Towards One-Way Functions from the Worst-case Hardness of Time-Bounded Kolmogorov Complexity.
\url{https://eccc.weizmann.ac.il/report/2025/125/}.
\textit{Further reference; not a proof endorsement.}
\item[S4]Capturing One-Way Functions via NP-Hardness of Meta-Complexity.
\url{https://eccc.weizmann.ac.il/report/2023/037/}.
\textit{Further reference; not a proof endorsement.}
\item[E1]Adi Akavia, Oded Goldreich, Shafi Goldwasser and Dana Moshkovitz,
\emph{On Basing One-Way Functions on NP-Hardness}, STOC 2006, 701--710,
author-hosted proceedings paper corresponding to [S1].
\url{https://people.csail.mit.edu/akavia/2006-stocAGGM.pdf}.
\textit{Reduction-model context; read together with the erratum below.}
\item[E2]Adi Akavia, Oded Goldreich, Shafi Goldwasser and Dana Moshkovitz,
\emph{Errata for: On Basing One-Way Functions on NP-Hardness},
8 February 2010.
\url{https://people.csail.mit.edu/akavia/AGGM_errata.pdf}.
\textit{Identifies the gap in the general adaptive argument and specifies
unaffected cases. Additional references consulted 27 September 2026.}
\end{itemize}
\end{document}
